%==============================================================================
\section{Act V: The Builder (\texttt{LocalMapping})}
%==============================================================================

While Tracking estimates poses, \texttt{LocalMapping} builds the map in the background.

\subsection{Why a Separate Thread?}

Local Bundle Adjustment (LBA) is too slow to run in real-time. If we did LBA for every frame, the system would lag behind. Instead:

\begin{itemize}
    \item \textbf{Tracking} runs fast (real-time) and just estimates pose
    \item \textbf{LocalMapping} runs slower but builds a better map
    \item They communicate through a queue of KeyFrames
\end{itemize}

\subsection{The Local Mapping Loop}

\begin{lstlisting}[caption={LocalMapping::Run() simplified}]
void LocalMapping::Run()
{
    while(!mbFinishRequested)
    {
        // Wait for new KeyFrame from Tracking
        if(CheckNewKeyFrames())
        {
            // 1. Process the new KeyFrame
            ProcessNewKeyFrame();
            
            // 2. Remove bad MapPoints
            MapPointCulling();
            
            // 3. Create new MapPoints by triangulation
            CreateNewMapPoints();
            
            // 4. Fuse duplicate MapPoints
            SearchInNeighbors();
            
            // 5. Local Bundle Adjustment
            if(!CheckNewKeyFrames())
            {
                Optimizer::LocalBundleAdjustment(mpCurrentKeyFrame, ...);
            }
            
            // 6. Remove redundant KeyFrames
            KeyFrameCulling();
        }
    }
}
\end{lstlisting}

\subsection{Step 1: ProcessNewKeyFrame}

When a new KeyFrame arrives:

\begin{lstlisting}[caption={ProcessNewKeyFrame}]
void LocalMapping::ProcessNewKeyFrame()
{
    // 1. Get the KeyFrame from the queue
    mpCurrentKeyFrame = mlNewKeyFrames.front();
    mlNewKeyFrames.pop_front();
    
    // 2. Compute Bag-of-Words
    mpCurrentKeyFrame->ComputeBoW();
    
    // 3. Associate MapPoints to the new KeyFrame
    for(int i = 0; i < mpCurrentKeyFrame->N; i++)
    {
        MapPoint* pMP = mpCurrentKeyFrame->mvpMapPoints[i];
        if(pMP)
            pMP->AddObservation(mpCurrentKeyFrame, i);
    }
    
    // 4. Update covisibility graph
    mpCurrentKeyFrame->UpdateConnections();
    
    // 5. Insert into the map
    mpAtlas->GetCurrentMap()->AddKeyFrame(mpCurrentKeyFrame);
}
\end{lstlisting}

\subsection{Step 3: CreateNewMapPoints}

This is where new 3D points are born:

\begin{lstlisting}[caption={CreateNewMapPoints simplified}]
void LocalMapping::CreateNewMapPoints()
{
    // For each neighboring KeyFrame
    for(KeyFrame* pKF2 : vpNeighKFs)
    {
        // Match features between current KF and neighbor
        vector<pair<int, int>> vMatchedIndices;
        ORBmatcher::SearchForTriangulation(mpCurrentKeyFrame, pKF2, 
                                           vMatchedIndices);
        
        // For each match, triangulate a new 3D point
        for(auto& match : vMatchedIndices)
        {
            cv::Point3f x3D = Triangulate(match);
            
            // Create new MapPoint
            MapPoint* pMP = new MapPoint(x3D, mpCurrentKeyFrame);
            pMP->AddObservation(mpCurrentKeyFrame, match.first);
            pMP->AddObservation(pKF2, match.second);
            
            mpAtlas->GetCurrentMap()->AddMapPoint(pMP);
        }
    }
}
\end{lstlisting}

\subsection{Step 5: Local Bundle Adjustment}

This is the expensive step that refines everything:

\begin{lstlisting}[caption={Local BA concept}]
void Optimizer::LocalBundleAdjustment(KeyFrame* pKF)
{
    // Collect local KeyFrames (connected to pKF)
    list<KeyFrame*> lLocalKeyFrames;
    
    // Collect local MapPoints (seen by local KFs)
    list<MapPoint*> lLocalMapPoints;
    
    // Collect fixed KeyFrames (neighbors of local KFs)
    list<KeyFrame*> lFixedKeyFrames;
    
    // Build g2o optimization graph
    g2o::SparseOptimizer optimizer;
    
    // Add pose vertices (local KFs)
    // Add point vertices (local MPs)
    // Add edges (reprojection errors)
    
    // Optimize!
    optimizer.optimize(5);  // 5 iterations
    
    // Update KeyFrame poses and MapPoint positions
}
\end{lstlisting}

\subsection{Step 6: KeyFrame Culling}

Remove redundant KeyFrames to keep the map compact:

\begin{lstlisting}[caption={KeyFrameCulling concept}]
void LocalMapping::KeyFrameCulling()
{
    // For each local KeyFrame
    for(KeyFrame* pKF : mvpLocalKeyFrames)
    {
        // Count how many of its MapPoints are also seen
        // by at least 3 other KeyFrames
        int nRedundantObs = 0;
        
        for(MapPoint* pMP : pKF->mvpMapPoints)
        {
            if(pMP->Observations() >= 3)
                nRedundantObs++;
        }
        
        // If 90% of points are redundant, cull the KeyFrame
        if(nRedundantObs > 0.9 * pKF->TrackedMapPoints())
        {
            pKF->SetBadFlag();  // Mark for deletion
        }
    }
}
\end{lstlisting}

\subsection{Communication with Tracking}

The interaction is through shared state:

\begin{lstlisting}[caption={Tracking to LocalMapping communication}]
// In Tracking::CreateNewKeyFrame()
mpLocalMapper->InsertKeyFrame(pKF);

// In LocalMapping
void InsertKeyFrame(KeyFrame* pKF)
{
    unique_lock<mutex> lock(mMutexNewKFs);
    mlNewKeyFrames.push_back(pKF);
}
\end{lstlisting}
